% !TEX program = xelatex
\documentclass[12pt,a4paper]{article}
\usepackage[T1]{fontenc}
%\usepackage{tgpagella, eulervm}

% --- XeLatex/LuaLatex ---
\usepackage[no-math]{fontspec}
\IfFontExistsTF{Fira Sans}{
	\setmainfont{Fira Sans}
}

% % --- This is for pdflatex ---
%\usepackage[light, sfdefault]{FiraSans}

\usepackage[italic]{mathastext}
\usepackage[bahasa]{babel}
\usepackage{booktabs}
\usepackage{geometry}
\usepackage{tabularx}
\usepackage{graphicx}
\usepackage{float}
\usepackage{xcolor}
\usepackage{hyperref}
\hypersetup{
	colorlinks=true,  % Colours links instead of boxes
	urlcolor=blue,    % Colour for external hyperlinks
	linkcolor=blue,   % Colour of internal links (sections, etc.)
	citecolor=red,    % Colour of citations
	bookmarks=true,   % Adds PDF bookmarks
	hypertexnames=true
}

\geometry{left=2cm, right=2cm, top=2cm}
\title{\textbf{Panduan Penggunaan Tangra}}
\author{Agus T. P. Jatmiko}
\hyphenation{me-la-ku-kan peng-amat-an deng-an}
\begin{document}
\maketitle
\noindent\rule{\linewidth}{0.4pt}
\section*{Instalasi}
\begin{enumerate}
	\item Unduh file instalasi Tangra dari laman Okultasi 2026 Observatorium Bosscha dan ekstrak ke folder kerja anda.
	\begin{center}
		\includegraphics[width=0.35\linewidth]{tangra-download}
	\end{center}
	\item Unduh file contoh data dari laman Okultasi 2026. Ekstrak pada folder kerja, berdampingan dengan folder Tangra untuk memudahkan.
	\begin{center}
		\includegraphics[width=0.2\linewidth]{tangra-data-download}
	\end{center}
	\item Klik ganda pada file \texttt{Tangra.exe} maka jendela utama Tangra akan muncul. Lakukan upgrade ke versi terbaru sebelum menggunakannya (klik pada tulisan dalam kotak merah). Tunggu sampai proses upgrade selesai.
	\begin{center}
		\includegraphics[width=0.5\linewidth]{tangra-upgrade}
	\end{center}
	\begin{center}
		\includegraphics[width=0.5\linewidth]{tangra-upgrade-progress}
	\end{center}
\end{enumerate}

\section*{Aplikasi Ekstraksi Kurva Cahaya}
\begin{enumerate}
	\item Klik \textit{File > Open Video} untuk membuka file contoh data (ekstensi .adv). Klik OK pada jendela yang muncul (biarkan pilihan default). Maka jendela utama akan muncul dan menampilkan data.
	\begin{center}
		\includegraphics[width=0.5\linewidth]{tangra-opendata}
	\end{center}
	\begin{center}
		\includegraphics[width=0.3\linewidth]{tangra-lc-reduction}
	\end{center}
	\underline{Catatan:}
	Jika tampilan data di jendela utama kurang jelas, atur rentang dinamik dengan \textit{slider} yang disediakan dengan mengaksesnya seperti gambar berikut. Tekan \textbf{Close} setelah yakin dengan hasilnya.
	\begin{center}
		\includegraphics[width=0.4\linewidth]{tangra-main-window}
	\end{center}
	\begin{center}
		\includegraphics[width=0.4\linewidth]{tangra-dynrange2}
	\end{center}
	 
	\underline{Catatan:} 
	\begin{itemize}
		\item Jika file data berupa *.fits, buka dengan \textit{File > Open FITS Sequence}. Arahkan ke folder contoh data (tekan bagian yang ditandai panah).
		\begin{center}
			\includegraphics[width=0.4\linewidth]{tangra-open-fits}
		\end{center}
		\item Jika data diambil dengan kamera CMOS, pilih \textbf{Video}, tekan OK.
		\begin{center}
			\includegraphics[width=0.4\linewidth]{tangra-open-fits-video}
		\end{center}
		\item Pilih \texttt{DATE-OBS} untuk memastikan \textit{timing} yang dipilih adalah waktu paparan di awal \textit{frame} dan \texttt{EXPTIME}.
		\begin{center}
			\includegraphics[width=0.4\linewidth]{tangra-open-fits-header}
		\end{center}
		\item Jika file data ada masalah di akurasi \textit{timing}, akan muncul jendela koreksi. Tekan OK.
		\begin{center}
			\includegraphics[width=0.4\linewidth]{tangra-open-fits-jitter}
		\end{center}
		\item Atur rentang dinamis dengan menggunakan \textit{slider} sampai obyek/bintang jelas terlihat dan tekan OK.
		\begin{center}
			\includegraphics[width=0.4\linewidth]{tangra-dynrange}
		\end{center}
		\item Sampai di sini, prosesnya mirip dengan proses membuka file .adv. Tekan OK pada jendela berikutnya, sampai file terbuka sempurna di Tangra.
		\begin{center}
			\includegraphics[width=0.4\linewidth]{tangra-lc-reduction}
		\end{center}
		\begin{center}
			\includegraphics[width=0.4\linewidth]{tangra-main-window}
		\end{center}
	\end{itemize}
	\item Klik pada bintang target dan klik \textbf{Add Object}.
	\begin{center}
		\includegraphics[width=0.4\linewidth]{tangra-add-object}
	\end{center}
	\item Pastikan pilihan seperti pada gambar dan tekan \textbf{Add}. Jika tidak ada bintang lain selain bintang target, pilih \textbf{Auto tracked}.
	\begin{center}
		\includegraphics[width=0.4\linewidth]{tangra-add-object2}
	\end{center}
	\item Lakukan untuk bintang--bintang yang lain sebagai bintang pembanding/\textit{guiding}. Pastikan \textit{setting} seperti pada gambar dan tekan \textbf{Add}.
	\begin{center}
		\includegraphics[width=0.4\linewidth]{tangra-add-comp}
	\end{center}
	\underline{Catatan:} Jika diperlukan, atur aperture untuk masing--masing bintang secara terpisah (bisa juga menggunakan aperture yang sama untuk semua bintang, sesuaikan dengan data yang dimiliki) dengan klik \textbf{Adjust}. Tekan OK jika sudah.
	\begin{center}
		\includegraphics[width=0.5\linewidth]{tangra-adjust}
	\end{center}
	\item Tekan \textbf{Start} untuk memulai ekstraksi data. Setelahnya, jendela baru akan muncul, menampilkan kurva cahaya masing--masing bintang.
	\begin{center}
		\includegraphics[width=0.5\linewidth]{tangra-start-extract}
	\end{center} 
	\begin{center}
		\includegraphics[width=0.5\linewidth]{tangra-hasil}
	\end{center} 
	\item Kurva cahaya yang ditampilkan bisa dikustomisasi dengan memilih bintang--bintang yang relevan (opsional).
	\begin{center}
		\includegraphics[width=0.5\linewidth]{tangra-hasil-pilih}
	\end{center}
	\underline{Catatan:}
	Pilih \textit{Data > Quick Re-process} untuk mencoba beberapa pilihan (lihat kotak merah) untuk mendapatkan hasil terbaik.
	\begin{center}
		\includegraphics[width=0.5\linewidth]{tangra-reprocess}
	\end{center}
	\begin{center}
		\includegraphics[width=0.5\linewidth]{tangra-reprocess-pilihan}
	\end{center}
	\item Jika sudah sesuai, export hasilnya ke dalam file .csv untuk pengolahan data lebih lanjut (dengan AOTA atau program lain) atau pilih \textit{File > Save Light Curve} untuk menyimpannya ke dalam file .lc untuk analisis dengan AOTA.
	\begin{center}
		\includegraphics[width=0.5\linewidth]{tangra-export}
	\end{center}
\end{enumerate}
\end{document}